% Part: lambda-calculus
% Chapter: lambda-definability
% Section: minimization

\documentclass[../../../include/open-logic-section]{subfiles}

\begin{document}

\olfileid{lam}{ldf}{min}
\olsection{લઘુત્તમીકરણ}

સામાન્ય પુનરાવર્તી વિધેયો એવાં વિધેયો છે, જે મૂળભૂત
વિધેયો $\Zero$, $\Succ$, $\Proj{n}{i}$માંથી સંયોજન,
આદિમ પુનરાવર્તન અને નિયમિત લઘુત્તમીકરણ વડે મેળવી
શકાય છે. બધા સામાન્ય પુનરાવર્તી વિધેયો
!!{lambda definable} છે એવું બતાવવા, !!a{lambda
  definable} વિધેય પર નિયમિત લઘુત્તમીકરણ લાગુ કરતાં
મળતું કોઈ પણ વિધેય પોતે !!{lambda definable} છે,
એ બતાવવું પડશે.

\begin{lem}
  \ollabel{lem:min} જો $f(x_1, \dots, x_k, y)$
  નિયમિત અને !!{lambda definable} હોય, તો
  \[
  g(x_1, \dots, x_k) = \umin{y}{f(x_1,\dots,x_k, y) = 0}
  \]
  વડે વ્યાખ્યાયિત $g$ પણ !!{lambda definable} છે.
\end{lem}

\begin{proof}
  ધારો કે લૅમ્બડા પદ~$F$ નિયમિત વિધેય
  $f(\vec x, y)$ને $\lambda$-વ્યાખ્યાયિત કરે છે.
  $g$ને !!{lambda define} કરવા આપણે શોધ-વિધેય
  અને સ્થિરબિંદુ સંયોજક વાપરીશું:
  \begin{align*}
    \fn{Search} & \ident \lambd[g][\lambd[f\,\vec{x}\,y][
        \fn{IsZero} (f\, \vec{x}\, y)\, y\, (g\, \vec{x} (\fn{Succ}\, y))]]\\
    H & \ident \lambd[\vec x][(Y \, \fn{Search}) F\, \vec{x}\, \num{0}],
  \end{align*}
  જ્યાં $Y$ કોઈ પણ સ્થિરબિંદુ સંયોજક છે.
  \sourcecorrection{OLLAM-060}{સ્થિર અંગ્રેજી મૂળની
  લઘુત્તમીકરણ-પ્રમેયમાં મેળવવાનું વિધેય~$g$ છે,
  પરંતુ પુરાવાની શરૂઆત અને અંતે ભૂલથી~$h$ આવે છે.
  અહીં વિધેય માટે~$g$ રાખ્યું છે; તેનું નિરૂપણ કરતું
  લૅમ્બડા પદ~$H$ મૂળ મુજબ જ રાખ્યું છે.}
  \sourcecorrection{OLLAM-061}{સ્થિર અંગ્રેજી મૂળમાં
  $\fn{Search}$ના પદમાં $(g\,\vec{x}(\fn{Succ}\,y))$
  આહ્વાનનું અંતિમ કૌંસ ગાયબ છે. અહીં તે ઉમેર્યું છે,
  જેથી $\fn{IsZero}$ની ત્રીજી શાખા યોગ્ય રીતે બંધ થાય.}
  અનૌપચારિક રીતે, $\fn{Search}$ પોતાનો સંદર્ભ લેતું
  વિધેય છે: $y$થી શરૂ કરીને $f\,\vec x\,y$
  શૂન્ય છે કે નહીં તે તપાસે છે; હોય તો $y$ પરત કરે
  છે, નહીંતર પોતાને $\fn{Succ}\,y$ સાથે ફરી
  બોલાવે છે. તેથી $(Y \, \fn{Search}) F
  \num{n_1}\dots\num{n_k}\,\num{0}$ એવું નાનામાં
  નાનું~$m$ પરત કરે છે કે જેના માટે $f(n_1,
  \dots, n_k, m) = 0$.

  ખાસ કરીને ધ્યાન આપો:
  \begin{align*}
    (Y \, \fn{Search}) F \num{n_1}
    \dots\num{n_k}\, \num{m} & \red \num{m}
    \intertext{જો $f(n_1, \dots,
      n_k, m) = 0$ હોય; અથવા}
    & \red (Y \, \fn{Search}) F\, \num{n_1} \dots\num{n_k}\, \num{m+1}
    \intertext{અન્યથા. $f$ નિયમિત હોવાથી કોઈક $y$ માટે
      $f(n_1, \dots, n_k, y) = 0$ થાય છે, તેથી}
    (Y \, \fn{Search}) F \num{n_1} \dots\num{n_k}\,\num{0}
    & \red \num{g(n_1, \dots, n_k)}.
    \end{align*}
\end{proof}


\begin{prop}
  દરેક સામાન્ય પુનરાવર્તી વિધેય !!{lambda definable} છે.
\end{prop}

\begin{proof}
 \olref[prf]{lem:basic} મુજબ બધા મૂળભૂત વિધેયો
 !!{lambda definable} છે; અને \olref[prf]{lem:comp},
 \olref[prf]{lem:prim} તથા \olref{lem:min} મુજબ
 !!{lambda definable} વિધેયોનો ગણ સંયોજન,
 આદિમ પુનરાવર્તન અને નિયમિત લઘુત્તમીકરણ હેઠળ
 બંધ છે.
\end{proof}

\end{document}
